\section{Introdução}

O câncer do colo do útero permanece entre as principais causas evitáveis de
morbidade e mortalidade por câncer em mulheres. Estimativas globais para 2020
indicam aproximadamente 604 mil novos casos e 342 mil mortes, com maior carga
em regiões que enfrentam restrições de acesso à vacinação, ao rastreamento e ao
tratamento oportuno \cite{singh2023global}. A citologia cervical continua sendo
um componente relevante do rastreamento, mas a leitura manual de lâminas é
trabalhosa, depende de especialistas e envolve alterações morfológicas que
podem ser sutis, especialmente nas categorias limítrofes.

Redes neurais profundas têm apresentado resultados promissores na análise de
imagens citológicas \cite{jiang2023systematic,xu2023cervical}. Entretanto,
revisões recentes apontam que a comparação
entre métodos é dificultada por diferenças de base de dados, preparação da
lâmina, tarefa, granularidade dos rótulos e protocolo de validação
\cite{fang2024systematic}. Um risco particularmente
importante surge quando células recortadas da mesma imagem aparecem em treino
e teste. Como elas compartilham coloração, aquisição, fundo e artefatos, uma
divisão aleatória no nível celular pode produzir uma estimativa excessivamente
otimista da generalização.

Além da discriminação, um sistema de apoio à triagem precisa tornar explícita a
relação entre falsos negativos e falsos positivos. Um limiar fixo de 0,5 não é
necessariamente apropriado para um cenário em que deixar de sinalizar uma
célula anormal pode ser mais grave do que encaminhar uma célula normal para
revisão. Probabilidades também não devem ser interpretadas como confiança sem
avaliação de calibração, pois redes modernas podem ser corretas na ordenação e,
ainda assim, superconfiantes \cite{guo2017calibration}. Diretrizes como
TRIPOD+AI reforçam a necessidade de descrever partições, seleção de limiar,
incerteza e uso pretendido de modelos preditivos \cite{collins2024tripod}.

Este artigo estuda a classificação binária \emph{normal versus anormal} de
células da coleção CRIC Cervix, obtidas por citologia convencional. A tarefa é
formulada como triagem: reunir ASC-US, LSIL, ASC-H, HSIL e SCC na classe
anormal, preservando os subtipos para análise posterior. Utiliza-se
ConvNeXt-Tiny pré-treinada, uma arquitetura convolucional modernizada que
incorpora decisões de projeto popularizadas por transformadores visuais
\cite{liu2022convnext}.

As contribuições do estudo são quatro. Primeiro, a avaliação primária emprega
validação cruzada em cinco folds, com agrupamento por imagem-fonte e retreino
do modelo em cada fold. Segundo, uma partição retida complementa a estimativa
de variância com análise de limiares operacionais, TTA, calibração e intervalos
de confiança por bootstrap agrupado por imagem. Terceiro, a ConvNeXt-Tiny é
comparada a uma ResNet50 treinada sob a mesma partição retida, evitando
comparação indireta entre bases incompatíveis. Quarto, os erros são
estratificados pelas categorias Bethesda originais, revelando onde a
simplificação binária permanece difícil.

O objetivo não é propor um diagnóstico autônomo nem comparar arquiteturas em
bases incompatíveis. Busca-se oferecer um protocolo reprodutível e
clinicamente orientado para avaliar até que ponto um classificador de células
pode priorizar regiões suspeitas em imagens convencionais, bem como explicitar
as evidências ainda necessárias antes de qualquer uso assistencial. Assim, o
diferencial do trabalho está menos em uma nova arquitetura e mais em uma
avaliação que combina controle de agrupamento, análise operacional de limiares,
incerteza estatística e interpretação por subtipo Bethesda em uma base de
citologia convencional.
